\section{Discussion}
\label{sec:disc}
Throughout our analysis of ASB 018 and five audit reports, we found that all five audits can be interpreted as compliant with ASB 018 despite falling short of the standard’s goals. We identified gaps between ASB 018’s envisioned practices and the practices it enables across five audit stages and highlighted how the design of the standard enables these gaps. 


First, we join \citet{berman2024scoping} in emphasizing the importance of tool developers providing clear definitions of the outcomes they intend for their tool to create. Our qualitative coding of ASB 018 and its accompanying factsheet to develop our interpretations of the goals of the tool’s creators was crucial to our evaluation of ASB 018 because the ASB did not explicitly state the standard’s goals. Without clear statements of intended outcomes, how can one conclude that a tool is effective? Furthermore, in situations like ours where tool developers do not clearly state goals, we contend that clearly articulating perceived goals from the outside, using methods such as our interpretive analysis of the tool’s goals, can help initiate conversations about desired outcomes and compel developers of audit tools to clarify their intended goals, towards holding tool creators (e.g., standards developers) accountable for the tools (e.g., standards) they create and promote. 

We additionally make the following recommendations, which we expand on in the sections below:
\begin{itemize}
    \item Designers of audit standards and other forms of audit tooling should carefully consider how designing a tool for compatibility with users' current perspectives and practices can undermine the tool's effectiveness. This is especially crucial for audit standards where ineffective standards can further legitimize existing practices. (Section \ref{sec:disc-compatibility})
    \item Evaluating an audit standard's effectiveness against fine-grained articulations of the standard's desired outcomes can help identify where increased specificity is crucial for achieving desired outcomes. Existing audit frameworks (e.g., \cite{ojewale2025towards, raji2020closing}) can help increase the specificity with which we articulate these desired outcomes. (Section \ref{sec:disc-design-artic})
    \item Components of audit standards where increased specificity may be crucial to achieving desired outcomes are: (i) actors and actor responsibilities, (ii) minimum activities, and (iii) intermediary steps. (Section \ref{sec:disc-design-specificity})
    \item Co-designing audit standards with groups beyond standards bodies and users of standards is crucial for counterbalancing users' desires for flexibility and can encourage more robust definitions of an audit standard's goals. (Section \ref{sec:disc-design-codesign})
\end{itemize} 


\subsection{Carefully consider how designing a tool for compatibility with tool users’ current perspectives and practices can undermine the tool’s effectiveness.}
\label{sec:disc-compatibility}
Prior work in human-computer interaction and algorithm auditing has called for the design of standards that more closely align with the needs and perspectives of those using the standard~\cite{schor2024mind, ojewale2025towards}. For instance, \citet{schor2024mind} highlight a gap between an algorithmic system transparency standard and system designers’ understandings of transparency. They call for standards bodies to more closely align standards with designers' current practices and understandings of transparency. These calls echo recent observations by \citet{berman2024scoping} that much of existing RAI tool development and evaluation prioritizes tool usability, potentially at the cost of effectiveness. 


While designing standards to align with users’ needs, workflows, and work contexts can be crucial for encouraging tool use and achieving intended outcomes~\cite{madaio2020co, holstein2019improving, rakova2021responsible}, \textbf{our findings caution that designing a standard to enable flexibility and compatibility with users’ current perspectives and practices can simultaneously undermine the standard’s effectiveness}. While vague language supports user discretion and improves the standard’s compatibility with labs’ existing approaches, it does not hold labs accountable for carrying out critical auditing practices. For instance, while ASB 018's requirement that labs  ``consider the impact of over- and underestimation of the number of contributors'' is compatible with labs' existing practices of observing whether LRs increase or decrease when analysts over- and underestimate the number of contributors to the mixture, it fails to hold labs accountable for measuring how inaccurate estimations of the number of contributors impacts system accuracy, sensitivity, specificity, and precision (Section~\ref{sec:scope-design}). 


We contend that \textbf{identifying how a standard’s compatibility may undermine its effectiveness is \textit{especially} critical in the context of standards for algorithm auditing, where ineffective standards not only may fail to improve practices, but at the same time may further legitimize  existing practices.} Three of the five audit reports we examined were completed before the publication of ASB 018 and one audit report published afterwards does not seem to engage with ASB 018, yet all four could be interpreted to comply with the standard’s requirements, suggesting that ASB 018 is compatible with existing lab practices: these labs would not need to significantly change their existing practices to comply with ASB 018’s requirements. In the context of ASB 018, this raises the question: Was the standard designed to improve practices or, as \cite{morrison2020vacuous} warn, was it designed to provide post hoc legitimacy to existing practices? In the next section, we discuss strategies that future efforts to design effective audit standards may leverage to balance flexibility and compatibility with effectiveness and accountability.


\subsection{Recommendations for designing audit standards that effectively support accountability}
\label{sec:disc-design}
\subsubsection{Fine-grained articulation of the audit standard’s desired outcomes and evaluating the standard’s effectiveness at achieving those outcomes can help identify areas where increased specificity is crucial for achieving desired outcomes.}
\label{sec:disc-design-artic}
Prior work on standards and algorithm auditing has framed increased specificity in standards as a double-edged sword: increased specificity can improve audit quality by mandating specific actions and outcomes, but can hinder auditor discretion that may be crucial to make context-appropriate, situational judgments~\cite{raji2022outsider, costanza2022audits, jackson2015standards}. Based on our findings, we suggest that \textbf{fine-grained articulation of outcomes desired by standards bodies and evaluating the standard’s effectiveness using these intended outcomes can help locate specific components where increased specificity is crucial for achieving the standard’s intended outcomes}. For instance, whereas ASB 018's goal for audits to establish boundaries on software use does not demand that labs use a specific ASB 018-mandated metric or threshold to define acceptable performance, it \textit{does} require that labs specify some definition of what constitutes acceptable performance and that these specifications are falsifiable (Section~\ref{sec:stds-design}) -- forms of specificity that ASB 018 crucially omits. Similarly, whereas ASB 018's goal for audits to test samples representative of the range of scenarios that the lab typically encounters or intends to analyze does not demand that all labs test up to five-person mixtures, it does require that labs specify a range that they and others could use to hold them accountable for.

In the context of audit tooling, we recommend using existing auditing frameworks (e.g., \cite{ojewale2025towards, raji2020closing, abebe2022adversarial, lam2024framework, radiya2023sociotechnical}) to define such goals. \textbf{Using \citet{ojewale2025towards}'s framework helped increase the specificity with which we articulated ASB 018’s goals and examined audit practices.} For instance, instead of articulating only the desired outcomes with respect to audit outputs (e.g., establish boundaries), we were able to assess audit reports against practices the standard’s envisions for each audit stage (e.g., scoping the audit, defining performance expectations, and audit documentation). \textbf{This also allowed us to address stages of the audit process beyond evaluation, such as standards identification and audit communication, that are crucial for creating accountability}~\cite{ojewale2025towards}. Echoing calls by \citet{ojewale2025towards} for tools that support all stages of the audit process, we argue that an audit standard’s goals should be articulated with sufficient granularity to address practices throughout the audit process. 


\subsubsection{Examine which actors, actor responsibilities, activities, and intermediary steps must be explicitly specified to accomplish the goals of the standard.}
\label{sec:disc-design-specificity}
Recent work studying algorithm audits in hiring has called for increased specificity in how rules for auditing (e.g., audit standards, requirements in audit legislation) define which algorithmic systems are considered in-scope~\cite{wright2024null, groves2024auditing} and who counts as an independent auditor~\cite{groves2024auditing}. Building on this work, \textbf{our findings suggest three components of audit standards where increased specificity may be crucial to ensuring audit standards produce standards developers' desired outcomes and help create accountability: (i) actors and actor responsibilities, (ii) bare minimum practices, and (iii) intermediary steps.} 


First, we found evidence that audit responsibilities were distributed between forensic labs and PGS developers in ways that undermined ASB 018’s goals (Section~\ref{sec:scope-gap1-target}). Yet ASB 018’s requirements do not acknowledge and engage with this reality. Many of the standard's requirements were directed at ``internal validation studies'' (as opposed to specific actors), and the few requirements assigned explicitly to labs contained vague language that allows a wide range of lab involvement in audits. This lack of specificity with respect to actors and actor responsibilities not only undermines the effectiveness of the standard, but also risks papering over distributions of auditing responsibilities that raise additional concerns such as lack of auditor independence. Recent work by \citet{groves2024auditing} similarly calls for increased specificity with respect to various actors involved in auditing and their specific roles and responsibilities, highlighting that this specificity can help ease the relational work required of auditors who must navigate tensions between software developers, software users, and regulators. 
Based on our findings, we add that assigning explicit roles and responsibilities to specific actors can not only ensure desired audit practices, but can also make visible the various actors and their roles in shaping audits in ways that open auditing practices to assessment a wider group of stakeholders.


Second, we found ASB 018’s use of vague requirements language such as ``consider'' and ``address'' consistently allowed for audit practices that fell short of the standard’s desired outcomes. For instance, the standard requires that labs ``address [...] precision'', and we see that labs analyzed a small subset of the test set used in the audit’s sensitivity and specificity experiment (Section~\ref{sec:performance-gap1}). In fact, this vague language allows for even more cursory practices. For instance, in PBSO's audit of STRmix v2.6.2, instead of measuring the precision of the software version being validated, the lab suggested that their previous validation of STRmix v2.4 and their sensitivity and specificity studies were sufficient. \textbf{We recommend that designers of audit standards carefully assess the bare minimum an audit could do to fulfill a requirement, and reflect on what specific practices they seek to require instead of relying on vague language.}


Lastly, we found that labs omitted crucial steps necessary to achieve ASB 018’s goals, and ASB 018 enabled this through its failure to require those preliminary steps. For instance, ASB 018's requirement that ``[audits] shall include [samples] that represent [...] the range of actual casework samples intended for analysis with the system at the laboratory'' depends on labs first determining and specifying the range of samples they intend to analyze with the software, but ASB 018 does not require that labs do so (Section~\ref{sec:scope-design}). These omissions also limited our ability to more thoroughly confirm study compliance (Section~\ref{sec:methods}). For instance, similar to how \citet{wright2024null}'s ``null compliance'' signals a lack of publicly available information needed to make a concrete determination of compliance, we found that labs did not provide a specific range of DNA samples that we could use to concretely assess the range of DNA samples tested in the audit. Consequently, following our permissive approach to assessing compliance, we simply accepted a lab's claim that their test set covers the lab's intended range. Requiring these preliminary steps can ensure  that those assessing audit compliance with a standard have crucial information needed to make unambiguous conclusions about audit compliance. \textbf{To design audit standards that effectively produce desired outcomes and support rigorous, outside assessments of compliance, we recommend that standards developers consider breaking requirements out into smaller units to ensure that the standard adequately accounts for crucial preliminary steps that other requirements rely on.}



\subsubsection{Engaging stakeholders beyond those who will be using the standard can help complement users’ perspectives and can uncover critical misunderstandings or disagreements over desired audit practices.}
\label{sec:disc-design-codesign}
As researchers and policymakers increasingly call for multistakeholder engagement in the design of audit standards~\cite{schor2024mind, aistandardshub} and RAI tooling more broadly, we argue that \textbf{it is crucial that co-design efforts engage groups beyond the standard’s intended users.} This can help ensure that users’ desires for flexibility and compatibility with practice do not undermine effectiveness (as discussed earlier in Section~\ref{sec:disc-compatibility}). 

Furthermore, \textbf{this broader engagement can help reveal how desired goals for audit practices may differ and conflict among individual members of standards-developing organizations (e.g., members of the ASB), users of standards (e.g., forensic laboratories), and various audit stakeholders (e.g., defense attorneys~\cite{jin2024beyond, abebe2022adversarial}, prosecutors, independent experts in forensic DNA profiling~\cite{krane2022using, butler2024dna})}. For instance, forensic laboratories may advocate for audit practices that make efficient use of laboratory resources,  molecular biology researchers may call for more experiments on large sample sizes and diverse samples, and defense attorneys may advocate for more thorough documentation practices. In the context of the U.S. criminal legal system, building on work by \citet{abebe2022adversarial} and \citet{jin2024beyond} that center the critical role of defense scrutiny of forensic software and call for adversarial audits~\cite{abebe2022adversarial}, we argue that it is especially crucial to engage defense attorneys' adversarial perspectives, as adversarialism is crucial to protecting defendants’ rights and the truth-seeking goals of the criminal legal system as a whole~\cite{abebe2022adversarial}. In other auditing contexts, engaging adversarial perspectives can not only help reveal standards developers' blindspots and make visible a broader range of concerns, but can also highlight where there is less consensus in audit scope, methods, and goals~\cite{geiger2024making, raji2022outsider, birhane2024ai}.


Frameworks that encourage fine-grained specification of desired outcomes can not only help standards bodies improve the specificity of their standards, towards creating effective audit standards, but can also help scaffold efforts to co-design standards with diverse stakeholders. Similar to how boundary objects for contested concepts like privacy~\cite{mulligan2016privacy} and fairness~\cite{mulligan2019thing} can support communication and collaboration across diverse groups of stakeholders, we suggest that using frameworks for auditing (e.g., \cite{abebe2022adversarial, ojewale2025towards, raji2020closing}) can facilitate better understanding of potential misunderstandings or disagreements over core concepts, practices, and desired outcomes. 
